\documentclass[12pt,a4paper]{article}

% ---------------------------------------------------------------
% Packages
% ---------------------------------------------------------------
\usepackage[utf8]{inputenc}
\usepackage[T1]{fontenc}
\usepackage[margin=1in]{geometry}
\usepackage{setspace}
\usepackage{titlesec}
\usepackage{graphicx}
\usepackage{amsmath,amssymb}
\usepackage{array,booktabs}
\usepackage{hyperref}
\usepackage{cite}
\usepackage{algorithm}
\usepackage{algpseudocode}
\usepackage{xcolor}
\usepackage{float}
\usepackage{caption,subcaption}
\usepackage{enumitem}
\usepackage{parskip}
\usepackage{fancyhdr}
\usepackage{abstract}
\usepackage{microtype}

% ---------------------------------------------------------------
% Hyperref
% ---------------------------------------------------------------
\hypersetup{
    colorlinks = true,
    linkcolor  = blue!70!black,
    citecolor  = green!50!black,
    urlcolor   = blue!70!black
}

% ---------------------------------------------------------------
% Header / Footer
% ---------------------------------------------------------------
\pagestyle{fancy}
\fancyhf{}
\lhead{\textit{Access2Education -- Personalised Learning via Ensemble ML}}
\cfoot{\thepage}

% ---------------------------------------------------------------
% Section formatting
% ---------------------------------------------------------------
\titleformat{\section}{\large\bfseries}{{\Roman{section}.}}{0.5em}{}
\titleformat{\subsection}{\normalsize\bfseries}{\Alph{subsection}.}{0.5em}{}

\doublespacing

% ===============================================================
\begin{document}

% ---------------------------------------------------------------
% Title Page
% ---------------------------------------------------------------
\begin{titlepage}
    \centering
    \vspace*{2cm}
    {\LARGE\bfseries Access2Education: A Full-Stack Adaptive Learning
    System Built on Ensemble Machine Learning, SM-2 Spaced
    Repetition, and Style-Aware Content Recommendation\par}
    \vspace{1.5cm}
    {\large Sambit Sharma\par}
    \vspace{0.3cm}
    {\normalsize Department of Artificial Intelligence and Data Science\par}
    \vspace{0.6cm}
    {\normalsize\textit{Under the Supervision of}\par}
    \vspace{0.2cm}
    {\large\bfseries Dr.\ Vishal Jain\par}
    \vspace{0.1cm}
    {\normalsize Department of Computer Science and Engineering\par}
    \vspace{0.6cm}
    {\normalsize \today\par}
    \vfill
    \rule{0.8\linewidth}{0.4pt}
\end{titlepage}

% ---------------------------------------------------------------
% Abstract
% ---------------------------------------------------------------
\begin{abstract}
\noindent
The single biggest inefficiency in online learning today is that
every student gets pushed the same reading list in the same order,
as if individual differences in thinking, memory, and problem-solving
style do not exist. This paper describes \textbf{Access2Education},
a system I built from the ground up to fix that problem. A new
student begins by completing a 25-question aptitude instrument that
measures five cognitive dimensions: logical reasoning, verbal
comprehension, numerical skill, memory retention, and sustained
attention. Each dimension is scored on a 0--100 scale. Those five
numbers feed into a soft-voting ensemble that combines
\texttt{GradientBoostingClassifier} (150 trees, learning rate 0.1,
max depth 4), \texttt{SVC} with an RBF kernel ($C=10$,
$\gamma=\texttt{scale}$, Platt scaling on), and
\texttt{LogisticRegression} ($C=5$, 1000 max iterations). This
ensemble, trained on a 10\,000-record synthetic dataset, reached a
cross-validated accuracy of \textbf{93.1\,\%} on the actual trained
model saved to disk (as recorded in \texttt{model\_info.json} on
10 April 2026). It classifies students into one of four learning
profiles: \emph{visual\_learner}, \emph{conceptual\_thinker},
\emph{practice\_based}, and \emph{step\_by\_step}. A content-based
recommendation engine then scores every item in a JSON library on
content type (up to 50 points), tag affinity (up to 30 points), and
adaptive difficulty (up to 20 points), returning the six highest-
scoring unfinished resources to the student dashboard. Beyond
recommendations, the platform also implements an SM-2-based spaced
repetition engine that schedules concept reviews using the
Ebbinghaus forgetting curve and assigns learning-style-specific
review formats (e.g., video clips for visual learners, mini exercises
for practice-based learners). The backend is Python 3.11 with
FastAPI and Motor for non-blocking MongoDB access; the frontend is
React 18.3 with Vite 5.2, Zustand 4.5, Axios 1.7, Recharts 2.12,
and Framer Motion 11. A client-side Demo Mode reimplements the full
API in the browser using \texttt{localStorage}, enabling zero-
infrastructure evaluation on Netlify. Expert-rated evaluation of the
top-six recommendations yielded a mean relevance score of 4.31 out of
5, versus 2.17 for random selection.
\end{abstract}

\noindent\textbf{Keywords:} adaptive learning, learning style
classification, ensemble machine learning, spaced repetition, SM-2,
content recommendation, FastAPI, React, MongoDB, Zustand.

\newpage
\tableofcontents
\newpage

% ===============================================================
\section{Introduction}
% ===============================================================

There is a specific frustration I had in mind when I started building
this project. When you study online, the platform has no idea
whether you are someone who needs to see a diagram before an
explanation makes sense, or someone who needs to hold a concept in
their head for a few minutes before they feel ready to touch the
code. Those are real differences, and they affect how quickly
you learn and how long you remember it. But virtually every online
learning tool I have used treats the problem as if it does not exist.

The goal of Access2Education is to make early-session personalisation
possible. Not after two weeks of interaction data, not after the
platform has seen you scroll past twenty videos---from the very first
login. To do that, the system asks you one question upfront: ``What
are you? How do you think?'' It answers that question not through
a self-report questionnaire where students often give socially
desirable answers, but through a 25-question aptitude instrument that
actually \emph{measures} five cognitive abilities. The scores feed
into a machine learning model that assigns a learning style category.
That category then drives every content decision the platform makes
for the next 30 days.

There are three things this system does that I had not seen combined
in a single open-source project before:

\begin{enumerate}[leftmargin=*, label=(\arabic*)]
  \item A \textbf{five-dimensional psychometric assessment}---five
    questions each for logical reasoning, verbal comprehension,
    numerical skill, memory retention, and sustained attention---that
    produces a score vector
    $(s_{\text{log}},\,s_{\text{ver}},\,s_{\text{num}},\,
    s_{\text{mem}},\,s_{\text{att}}) \in [0,100]^5$.
  \item A \textbf{three-tier ML inference cascade}:
    soft-voting ensemble as the primary classifier (93.1\,\% CV
    accuracy on the saved model) $\rightarrow$ K-Means clustering
    fallback (92.0\,\%) $\rightarrow$ dominant-score rule mapping
    (safe default), so a style label is always returned regardless of
    disk state.
  \item A \textbf{learning-style-adapted SM-2 spaced repetition
    engine} that schedules concept reviews according to the
    Ebbinghaus forgetting curve and delivers review material
    in the student's preferred format (video for visual learners,
    mini-exercise for practice-based, etc.).
\end{enumerate}

The remainder of this paper is structured as follows. Section II
discusses related work. Section III describes the system architecture.
Section IV explains the ML pipeline in detail. Section V covers the
content recommendation engine. Section VI describes the spaced
repetition system. Section VII goes through the full-stack
implementation. Section VIII presents evaluation results. Section IX
discusses limitations and future directions. Section X concludes.

% ===============================================================
\section{Related Work}
% ===============================================================

\subsection{Adaptive Educational Systems}

Brusilovsky's foundational work on adaptive hypermedia
\cite{brusilovsky1996} demonstrated that altering what content a
learner sees based on a stored model of their knowledge state could
reduce task-completion time relative to static systems. Later
commercial platforms built on Bayesian Knowledge Tracing
\cite{corbett1994} scaled this idea to millions of students.
The persistent challenge in all of these systems is the observation
window before personalisation kicks in: a brand-new user on any
interaction-log-based platform is essentially un-personalised for
the first few sessions.

Access2Education bypasses this entirely by using a structured
aptitude instrument at registration. The style label is available
immediately, so the first page of recommendations a student sees is
already filtered for their cognitive profile.

\subsection{Learning Style Classification with Machine Learning}

Abdallah et al.\ \cite{abdallah2020} applied a Random Forest to
Felder-Silverman questionnaire responses, reaching about 89\,\%
accuracy. Truong \cite{truong2016} found that SVMs with RBF kernels
consistently outperform tree models and Naive Bayes on four-class
learning style tasks, reporting 92\,\% accuracy---which is consistent
with what I observed in my own model comparisons. Nafea et al.\
\cite{nafea2021} used K-Means to segment students in an Egyptian MOOC
and showed that cluster-matched content lifted quiz scores by 14\,\%.

My approach differs in one key way: rather than choosing a single
winning model, I combine three classifiers through soft-voting.
The reason is not to chase a higher mean accuracy number---individual
models are already in the 93\,\% range---but to reduce the variance
of predictions at the tails, which matters more in a deployment where
a wrong classification sticks for 30 days.

\subsection{Spaced Repetition}

The SM-2 algorithm \cite{supermemo} was designed by Piotr
Wo\'zniak for the SuperMemo project and forms the basis of most
modern flashcard systems (Anki being the most widely known). It
updates an easiness factor $EF$ and an inter-repetition interval $I$
after each review based on a quality rating from 0 (blackout) to 5
(perfect). What I added on top of the standard SM-2 is a
learning-style-specific review format selector: when a card is due,
the system does not just surface a generic question---it wraps the
question in the media format most likely to trigger recall for that
student's learning style.

\subsection{Content-Based Recommendation}

Huang et al.\ \cite{huang2019} argued that content-based filtering
outperforms collaborative filtering when item metadata is rich and
when users are new (no interaction history). That framing fits
exactly: the Access2Education recommendation engine scores items on
metadata only (type, tags, difficulty) and never needs a
user-to-user similarity matrix.

% ===============================================================
\section{System Architecture}
% ===============================================================

\subsection{Three-Tier Structure}

The platform is split into three layers with clean HTTP interfaces
between them.

\begin{enumerate}[leftmargin=*,label=\textbf{Tier \arabic*:}]
  \item \textbf{Frontend} --- React 18.3 SPA bundled by Vite 5.2.
    Zustand 4.5 manages global state. Framer Motion 11 handles
    page transitions. Recharts 2.12 and react-circular-progressbar
    render analytics. React Hook Form with Zod validation handles
    all forms. react-markdown with remark-gfm renders chatbot
    responses.
  \item \textbf{Backend} --- FastAPI on Python 3.11 with Motor
    (async MongoDB driver). Four router groups:
    \texttt{/auth}, \texttt{/test}, \texttt{/content},
    \texttt{/chatbot}. Pydantic v2 validates all requests and
    responses.
  \item \textbf{Data and Models} --- MongoDB 7.0 with five
    collections: \texttt{users}, \texttt{test\_results},
    \texttt{progress}, \texttt{refresh\_tokens}, \texttt{otps}.
    ML artefacts in \texttt{ML/models/}: \texttt{classifier.pkl}
    (1.5 MB), \texttt{scaler.pkl}, \texttt{label\_encoder.pkl},
    \texttt{kmeans\_model.pkl}, \texttt{cluster\_map.pkl}, and
    \texttt{model\_info.json}.
\end{enumerate}

\begin{figure}[H]
  \centering
  \begin{verbatim}
  Browser  (React 18 + Vite 5)
        |       HTTPS / REST
        v
  FastAPI  (Python 3.11 + Motor)
       |               |
       v               v
   MongoDB 7.0     ML/models/
   (5 collections)  (.pkl files)
  \end{verbatim}
  \caption{Simplified architecture of Access2Education.}
\end{figure}

\subsection{Authentication}

Authentication uses a dual-token JWT scheme implemented in
\texttt{backend/routes/auth.py}. A short-lived access token
(30-minute TTL, signed with \texttt{SECRET\_KEY} using HS256)
travels in the \texttt{Authorization: Bearer} header. A refresh
token (7-day TTL, signed with a separate \texttt{REFRESH\_SECRET\_KEY})
is stored in the \texttt{refresh\_tokens} MongoDB collection under
a TTL index on the \texttt{expires\_at} field.

Passwords are hashed using passlib's bcrypt context. The
\texttt{RegisterRequest} Pydantic model enforces a minimum length
of 8 characters, requires at least one uppercase letter and one
digit (validated with \texttt{@field\_validator}), and accepts only
the roles \texttt{student} or \texttt{teacher}.

On the frontend, an Axios 1.7 response interceptor handles
silent token refresh. When a 401 arrives, the interceptor sets a
shared \texttt{isRefreshing} flag before doing anything else. All
in-flight requests attach to a pending promise rather than each
independently firing a refresh call---which would create a race where
the server invalidates each token before the previous request
finishes. Once the refresh resolves, every queued request replays
with the new token.

Password recovery uses a six-digit OTP (random digits via
\texttt{secrets.token\_hex} equivalent, valid for 10 minutes,
TTL-indexed in the \texttt{otps} collection). The OTP is sent
via FastAPI \texttt{BackgroundTasks} so the HTTP response returns
immediately without waiting for the email dispatch.

\subsection{Demo Mode}

The Netlify deployment cannot run Python server-side. To support
it, I wrote \texttt{demoApi.js}: a single JavaScript file that
reimplements the full backend REST contract using
\texttt{localStorage}. Vite's module aliasing swaps every
\texttt{api.get()} and \texttt{api.post()} call for its mock
counterpart when \texttt{VITE\_DEMO\_MODE=true} is set at build time.

The mock implements: registration and login with
\texttt{localStorage}-backed user store, the full 25-question
aptitude flow, the style classification heuristic, content browsing
with pagination, progress tracking, and a keyword-matching chatbot
template engine. No network request is made. This let me ship a
fully interactive pilot on Netlify without any server, which turned
out to be genuinely valuable for demonstrations.

% ===============================================================
\section{Machine Learning Pipeline}
% ===============================================================

\subsection{Dataset}

Training data is handled by \texttt{ML/train\_cluster.py}. The
script first looks for a real CSV at
\texttt{data/student\_aptitude\_dataset.csv}. If the file is not
there (or is missing required columns), the function
\texttt{generate\_synthetic\_data(n=10000)} is called instead.
That function creates 10\,000 records distributed across four
class profiles using clipped Gaussians:

\begin{equation}
  s_i \;\sim\; \mathcal{N}(\mu_i, \sigma_i^2)
  \quad\text{clipped to}\;[0,100]
\end{equation}

Table 1 shows the class-specific means and standard deviations
taken directly from the \texttt{profiles} dictionary in
\texttt{generate\_synthetic\_data()}:

\begin{table}[H]
  \centering
  \caption{Synthetic dataset parameters: class distribution and
           Gaussian parameters
           $(\mu/\sigma)$ per cognitive dimension.}
  \small
  \begin{tabular}{lcccccc}
    \toprule
    \textbf{Class} & \textbf{N} &
    \textbf{Log.} & \textbf{Ver.} & \textbf{Num.} &
    \textbf{Mem.} & \textbf{Att.} \\
    \midrule
    visual\_learner      & 2,600 & 55/10 & 80/8  & 45/12 & 78/9  & 70/10 \\
    conceptual\_thinker  & 2,200 & 85/7  & 65/10 & 70/9  & 55/11 & 60/10 \\
    practice\_based      & 2,800 & 65/9  & 50/11 & 85/7  & 60/10 & 80/8  \\
    step\_by\_step       & 2,400 & 60/10 & 62/10 & 58/9  & 85/7  & 75/9  \\
    \midrule
    \textbf{Total} & \textbf{10,000} & & & & & \\
    \bottomrule
  \end{tabular}
\end{table}

After generation, the DataFrame is shuffled with
\texttt{random\_state=42} and \texttt{reset\_index(drop=True)}
before any splitting, ensuring reproducibility.

\subsection{Feature Standardisation}

The \texttt{preprocess()} function in \texttt{train\_cluster.py}
applies \texttt{StandardScaler} to the raw 0--100 scores:

\begin{equation}
  \tilde{x}_{ij} \;=\; \frac{x_{ij} - \bar{x}_j}{\sigma_j}
\end{equation}

The scaler is fitted on the full training set and serialised to
\texttt{ML/models/scaler.pkl} alongside the classifier, so the
exact same transformation is applied at inference time in
\texttt{predict\_cluster.py}: \texttt{\_SCALER.transform(features\_arr)}.
This is important because the SVM's RBF kernel is sensitive to
feature scale---without standardisation, a feature with a wider
raw distribution would dominate the kernel distance computation.

\subsection{Learning Style Operationalisation}

The four categories used in the system are grounded in Fleming's
VARK model \cite{fleming1992} and Kolb's experiential learning
cycle \cite{kolb1984}, but expressed through measurable aptitude
scores rather than self-report items. The operationalisation is
defined in \texttt{ML/recommender.py} through the
\texttt{LEARNING\_STYLE\_RULES} dictionary:

\begin{itemize}[leftmargin=*]
  \item \textbf{visual\_learner} --- Elevated verbal (mean 80) and
    memory (mean 78) scores. Preferred content types:
    \texttt{video}, \texttt{infographic}. Tag boosts:
    \emph{visual}, \emph{animation}, \emph{diagram}.
  \item \textbf{conceptual\_thinker} --- High logical score (mean
    85). Preferred types: \texttt{article}, \texttt{video}. Tag
    boosts: \emph{conceptual}, \emph{theory}, \emph{case-study}.
  \item \textbf{practice\_based} --- High numerical (mean 85) and
    attention (mean 80) scores. Preferred types:
    \texttt{exercise}, \texttt{project}. Tag boosts:
    \emph{practice}, \emph{hands-on}, \emph{coding}.
  \item \textbf{step\_by\_step} --- High memory (mean 85) with
    balanced attention (mean 75). Preferred types:
    \texttt{notes}, \texttt{tutorial}. Tag boosts:
    \emph{structured}, \emph{step-by-step}, \emph{memory}.
\end{itemize}

\subsection{Ensemble Classifier}

The \texttt{train\_ensemble()} function builds a
\texttt{VotingClassifier} with \texttt{voting='soft'} over three
estimators. Soft voting averages the posterior probability vectors
from each model before taking the argmax:

\begin{align}
  \hat{y} &= \arg\max_{k \in \{0,1,2,3\}}
  \frac{1}{3}\Bigl[
    \hat{p}_{\text{GB}}(k \mid \mathbf{x})
  + \hat{p}_{\text{SVM}}(k \mid \mathbf{x})
  + \hat{p}_{\text{LR}}(k \mid \mathbf{x})
  \Bigr]
\end{align}

\noindent\textbf{Gradient Boosting (GB).}
\texttt{n\_estimators=150}, \texttt{learning\_rate=0.1},
\texttt{max\_depth=4}, \texttt{random\_state=42}.
Gradient boosting fits each successive tree to the negative
gradient of the cross-entropy loss of the current ensemble, thereby
capturing feature interaction effects that a linear model misses.
The moderate depth of 4 keeps individual trees weak enough to
benefit from boosting without overfitting on the 10\,000-record
dataset.

\noindent\textbf{Support Vector Machine (SVM).}
\texttt{kernel='rbf'}, $C=10$, \texttt{gamma='scale'},
\texttt{probability=True}, \texttt{random\_state=42}.
The RBF kernel implicitly maps the standardised five-dimensional
input into a higher-dimensional space where the four class clusters
become linearly separable. Setting \texttt{probability=True}
enables Platt scaling \cite{platt1999}, which converts the SVM's
margin-based decision scores into calibrated class probabilities
required for meaningful soft voting.

\noindent\textbf{Logistic Regression (LR).}
$C=5$, \texttt{max\_iter=1000}, \texttt{random\_state=42}.
LR provides a well-calibrated linear baseline. Its probabilities
are less likely to be overconfident in low-data regions of the
feature space, which anchors the ensemble against situations where
the two non-linear models overfit to the same distributional
artefact.

\subsection{Cross-Validation and Model Comparison}

\texttt{train\_ensemble()} runs \texttt{StratifiedKFold(n\_splits=5,
shuffle=True, random\_state=42)} and evaluates accuracy with
\texttt{cross\_val\_score}. Table 2 shows all models I tested.

\begin{table}[H]
  \centering
  \caption{Five-fold stratified CV accuracy comparison.
           The ``Actual CV'' column for the ensemble reflects
           the value stored in \texttt{model\_info.json} from the
           last training run on 10 April 2026.}
  \begin{tabular}{lcc}
    \toprule
    \textbf{Model} & \textbf{Mean Accuracy} & \textbf{Std Dev} \\
    \midrule
    Voting Ensemble (GB+SVM+LR) & \textbf{93.1\,\%} & $\pm$0.6\,\% \\
    SVM (RBF, $C=10$)           & 93.4\,\%           & $\pm$0.3\,\% \\
    Gradient Boosting           & 93.2\,\%           & $\pm$0.3\,\% \\
    Logistic Regression         & 93.1\,\%           & $\pm$0.3\,\% \\
    K-Nearest Neighbours ($k=5$) & 92.2\,\%          & $\pm$0.5\,\% \\
    K-Means (unsupervised)       & 92.0\,\%          & ---           \\
    \bottomrule
  \end{tabular}
\end{table}

The standalone SVM has a marginally higher mean (93.4 vs 93.1).
I chose the ensemble for its worst-case stability argument: in any
single fold where one model makes a poor boundary decision, the
other two models outvote it, reducing the probability of a bad
classification at any particular point in the student feature
space. In a deployment where a misclassification drives the wrong
content stream for 30 days (the hardcoded re-test cooldown in
\texttt{test.py} as \texttt{RETEST\_COOLDOWN\_DAYS}), that
structural redundancy is worth more than 0.3\,\% mean accuracy.

\subsection{Confidence Score and Three-Tier Fallback}

The confidence score returned with every prediction is the maximum
averaged posterior across classes:

\begin{equation}
  c \;=\; \max_{k}\;\frac{1}{3}\Bigl[
    \hat{p}_{\text{GB}}(k \mid \mathbf{x})
  + \hat{p}_{\text{SVM}}(k \mid \mathbf{x})
  + \hat{p}_{\text{LR}}(k \mid \mathbf{x})
  \Bigr]
\end{equation}

This is stored in \texttt{test\_results.confidence} and shown on
the result page. The inference logic in
\texttt{ML/predict\_cluster.py} determines which of three methods
to use at startup by checking which model files exist on disk:

\begin{algorithm}[H]
\caption{Three-Tier Prediction in \texttt{predict\_learning\_style()}}
\begin{algorithmic}[1]
\If{\texttt{\_CLASSIFIER} and \texttt{\_SCALER} and
    \texttt{\_LABEL\_ENCODER} are not None}
  \State $\tilde{\mathbf{x}} \leftarrow
         \texttt{\_SCALER.transform}([\mathbf{x}])$
  \State $idx \leftarrow \texttt{\_CLASSIFIER.predict}(\tilde{\mathbf{x}})[0]$
  \State $\text{style} \leftarrow \texttt{\_LABEL\_ENCODER.classes\_}[idx]$
  \State $\text{proba} \leftarrow \texttt{\_CLASSIFIER.predict\_proba}(\tilde{\mathbf{x}})[0]$
  \State \textbf{return} $(\text{style},\; \text{proba}[idx])$
\ElsIf{\texttt{\_KMEANS} and \texttt{\_SCALER} and
       \texttt{\_CLUSTER\_MAP} are not None}
  \State $\tilde{\mathbf{x}} \leftarrow
         \texttt{\_SCALER.transform}([\mathbf{x}])$
  \State $k \leftarrow \texttt{\_KMEANS.predict}(\tilde{\mathbf{x}})[0]$
  \State \textbf{return} $(\texttt{\_CLUSTER\_MAP}[k],\; 0.85)$
\Else
  \State $d \leftarrow \arg\max_{j \in \text{FEATURES}}\; x_j$
  \State \textbf{return} $(\texttt{\_RULE\_BASED\_MAP}[d],\; 0.70)$
\EndIf
\end{algorithmic}
\end{algorithm}

The \texttt{\_mode} variable is set at import time and printed to
the console so it is always clear which tier is active.

\subsection{K-Means as Fallback}

The K-Means model is trained in \texttt{train\_kmeans()} with
\texttt{n\_clusters=4}, \texttt{n\_init=20}, \texttt{max\_iter=500}.
Since K-Means assigns arbitrary integer labels, a semantic
cluster-to-style mapping is computed using
\texttt{scipy.optimize.linear\_sum\_assignment}---the Hungarian
algorithm \cite{kuhn1955}---applied to the confusion matrix between
ground-truth labels and cluster assignments:

\begin{equation}
  \pi^* = \arg\min_\pi \sum_{i} -C_{i,\,\pi(i)}
\end{equation}

The resulting \texttt{cluster\_map} dictionary (e.g.
\texttt{\{0: 'visual\_learner', 2: 'step\_by\_step', ...\}})
is serialised to \texttt{cluster\_map.pkl} and loaded at inference
time.

% ===============================================================
\section{Content Recommendation Engine}
% ===============================================================

\subsection{Scoring Function}

The \texttt{Recommender} class in \texttt{ML/recommender.py} loads
the content library from \texttt{data/content\_metadata.json} at
initialisation. For each item not in the student's completed set,
\texttt{\_score\_item()} computes a relevance score out of 100
across three components:

\begin{equation}
  \text{score}(r_j) =
    \underbrace{S_{\text{type}}}_{0,25,\text{ or }50}
  + \underbrace{\min\!\bigl(30,\;|\text{tags}(r_j)
    \cap T_s^{\text{tags}}|\times 10\bigr)}_{\text{tag affinity}}
  + \underbrace{S_{\text{diff}}(n_c)}_{0,5,10,15,\text{ or }20}
\end{equation}

For the type component, a preferred type scores 50 points and a
secondary type scores 25 points (0 for no match). The difficulty
component is more granular. Looking directly at the code in
\texttt{\_score\_item()}:

\begin{table}[H]
  \centering
  \caption{Difficulty score $S_{\text{diff}}$ as a function of
           item difficulty level and student completion count $n_c$.}
  \begin{tabular}{lccccc}
    \toprule
    & \multicolumn{4}{c}{\textbf{Item difficulty}} \\
    \cmidrule(lr){2-5}
    \textbf{Completion count $n_c$} & \textbf{1} & \textbf{2} & \textbf{3} & \textbf{$\ge$3} \\
    \midrule
    $n_c = 0$    & 20 & 10 & 0  & 0 \\
    $n_c < 3$    & 20 & 15 & 0  & 0 \\
    $3 \le n_c < 8$ & 5  & 20 & 15 & 0 \\
    $n_c \ge 8$  & 0  & 10 & 20 & 20 \\
    \bottomrule
  \end{tabular}
\end{table}

This scheme nudges students toward progressively harder material as
their completion count grows, implementing a lightweight version of
Vygotsky's zone of proximal development \cite{vygotsky1978} without
requiring a formal IRT model.

\subsection{Sorting and Prerequisite Handling}

The final ranking is not purely by score. The \texttt{get\_recommendations()}
method sorts by the tuple \texttt{(prerequisites\_met, recommendation\_score)}
in descending order:

\[
  \text{sorted} = \text{sort}
  \Bigl[\bigl(\mathbf{1}[\text{prereqs met}],\;
  \text{score}\bigr)\Bigr]_{\text{desc}}
\]

This guarantees that an item with all prerequisites satisfied always
appears above a higher-scoring item that has an unmet prerequisite.
The check in \texttt{\_prerequisites\_met()} simply verifies that
every ID listed in \texttt{item["prerequisites"]} appears in
\texttt{completed\_ids}.

\subsection{Additional Endpoints}

The \texttt{Recommender} class also exposes:
\begin{itemize}[leftmargin=*]
  \item \texttt{search(query)}: substring match across title,
    subject, description, and tags for the \texttt{/content/search}
    route.
  \item \texttt{get\_all\_subjects()}: aggregates the library by
    subject, returning total item count and available types per
    subject.
  \item \texttt{get\_score\_matrix()}: returns the full
    style $\times$ item relevance matrix, useful for offline
    analysis.
\end{itemize}

% ===============================================================
\section{Spaced Repetition Engine}
% ===============================================================

\subsection{SM-2 Algorithm Implementation}

\texttt{ML/spaced\_repetition.py} implements the SM-2 algorithm
\cite{supermemo} within an \texttt{SM2Card} class. Each card stores
the standard SM-2 state fields: repetition count, easiness factor
($EF$, initialised at 2.5 and clamped to a minimum of 1.3), review
interval in days, total review count, and consecutive-correct
streak. The \texttt{review(quality)} method processes a quality
rating $q \in [0,5]$ and updates the state:

\begin{equation}
  EF' \;=\; \max\!\bigl(1.3,\;
  EF + 0.1 - (5-q)\cdot(0.08 + (5-q)\cdot 0.02)\bigr)
\end{equation}

If $q < 3$ (failed recall), the repetition counter and interval are
reset to zero and one respectively, and the streak resets. For
successful recalls ($q \ge 3$), the interval follows the standard
SM-2 schedule:

\begin{equation}
  I_n =
  \begin{cases}
    1  & \text{if } n = 1 \\
    6  & \text{if } n = 2 \\
    \lceil I_{n-1} \cdot EF' \rceil & \text{if } n > 2
  \end{cases}
\end{equation}

The next review date is then set to
$\texttt{datetime.utcnow()} + \texttt{timedelta(days=}I_n\texttt{)}$.

\subsection{Retention Prediction}

The \texttt{predict\_retention(days\_from\_now)} method estimates
the probability of successful recall at a future point using the
Ebbinghaus forgetting curve:

\begin{equation}
  R(t) \;=\; e^{-t/S}
\end{equation}

where $t$ is the number of days since the last review plus
\texttt{days\_from\_now}, and $S = \max(1,\; I \times 0.7)$ is a
stability estimate derived from the current interval. This is used
in two places: the \texttt{get\_upcoming\_schedule()} method
annotates each future card with its predicted retention at review
time, and \texttt{get\_retention\_heatmap()} aggregates per-subject
average retention for the dashboard visualisation.

\subsection{Urgency Score and Review Prioritisation}

When multiple cards are due, the \texttt{SpacedRepetitionEngine.
get\_due\_cards()} method sorts them by an urgency score:

\begin{equation}
  U = \underbrace{(1-R)\cdot 2}_{\text{retention penalty}}
  + \underbrace{\ln(1 + \max(0, -\text{days\_until\_due}))\cdot 0.5}
              _{\text{overdue bonus}}
  + \underbrace{\max(0,\,(2.5-EF)\cdot 0.3)}_{\text{difficulty factor}}
\end{equation}

Cards with low current retention are penalised heavily. Cards that
are overdue get a logarithmic bonus (so one very overdue card does
not completely dominate the queue). Cards with a low easiness factor
(historically hard to recall) receive an additional boost.

\subsection{Learning-Style-Specific Review Formats}

The \texttt{LEARNING\_STYLE\_REVIEW\_FORMAT} dictionary maps each
of the four style indices to a specific review media type:

\begin{table}[H]
  \centering
  \caption{Learning-style-specific review formats in the SM-2 engine.}
  \begin{tabular}{lll}
    \toprule
    \textbf{Style} & \textbf{Primary format} & \textbf{Prompt prefix} \\
    \midrule
    Visual Learner     & video\_clip       & ``Watch this 2-min visual recap, then answer:'' \\
    Conceptual Thinker & article\_excerpt  & ``Read this conceptual overview, then answer:'' \\
    Practice-Based     & mini\_exercise    & ``Solve this quick problem to test recall:'' \\
    Step-by-Step       & step\_tutorial    & ``Review these steps, then answer:'' \\
    \bottomrule
  \end{tabular}
\end{table}

This goes beyond standard SM-2 by ensuring that the review
experience itself is calibrated to how the student most easily
retrieves information, not just \emph{when} they review.

% ===============================================================
\section{Full-Stack Implementation}
% ===============================================================

\subsection{Backend Routes}

The backend exposes six route groups. The four primary ones
relevant to the ML pipeline are described here.

\noindent\textbf{\texttt{/auth} (9 endpoints).}
Register, login, OAuth2 form login (for Swagger UI), token refresh,
logout, get profile, update profile (name, avatar, password),
forgot-password (OTP dispatch via \texttt{BackgroundTasks}), and
reset-password. The \texttt{get\_current\_student} dependency chain
validates the access token and additionally enforces
\texttt{role == "student"} before any test or content route proceeds.

\noindent\textbf{\texttt{/test} (6 endpoints).}
Fetch questions, submit test, get latest result, get history,
retake check, and model info. The submit route (\texttt{POST
/test/submit}) is the core of the system. It:
\begin{enumerate}[leftmargin=*, nosep]
  \item queries \texttt{test\_results} for the most recent submission;
    returns HTTP 429 if within the 30-day cooldown;
  \item calls \texttt{calculate\_section\_scores()} to score all 25
    answers against \texttt{QUESTIONS\_BANK} and compute section-wise
    0--100 values;
  \item calls \texttt{run\_ml\_prediction()} which invokes
    \texttt{predict\_cluster\_id()} and \texttt{predict\_learning\_style()}
    from \texttt{ML/predict\_cluster.py};
  \item inserts the result document (including \texttt{ml\_mode}:
    \texttt{"ensemble"} or \texttt{"rule\_based"}) into
    \texttt{test\_results};
  \item updates the \texttt{users} document with \texttt{learning\_style},
    \texttt{cluster\_id}, \texttt{last\_test\_date}, and
    \texttt{needs\_recluster: False};
  \item returns scores, style, confidence, and
    \texttt{next\_test\_date}.
\end{enumerate}

\noindent\textbf{\texttt{/content} (8 endpoints).}
Personalised recommendations, all-content browse with pagination
(subject/type/difficulty filters, page and limit query params),
subjects list, full-text search, progress summary with
subject-level breakdown, content detail with aggregated rating,
mark-complete (upsert into \texttt{progress}), and rate-content.

\noindent\textbf{\texttt{/chatbot} (1 endpoint).}
Proxies the DeepSeek LLM API with two prompt templates:
summarisation (condenses pasted lecture text to bullet points) and
doubt resolution (answers a student question with optional lecture
context). Demo Mode replaces this with a keyword template function.

\subsection{Frontend Stack}

The frontend depends are versioned exactly in
\texttt{Frontend/package.json}: React 18.3.1, react-router-dom
6.23.1, Axios 1.7.2, Recharts 2.12.7, react-circular-progressbar
2.1.0, Framer Motion 11.2.4, Zustand 4.5.2, react-hook-form 7.51.5
with Zod 3.23.8 validation, react-markdown 9.0.1 with remark-gfm,
lucide-react 0.379.0 for icons, react-hot-toast 2.4.1, and dayjs
1.11.11 for date formatting. The build tool is Vite 5.2.11 with
\texttt{@vitejs/plugin-react} 4.3.0 and Tailwind CSS 3.4.3.

Zustand's \texttt{persist} middleware backs the auth slice to
\texttt{localStorage}, restoring the user session silently on page
refresh without a round-trip to the server. The progress store
updates reactively via Zustand selectors, so Recharts charts re-render
automatically whenever a content-completion event mutates the store.

\subsection{Deployment}

\noindent\textbf{Docker Compose.} A \texttt{docker-compose.yml}
orchestrates three services: frontend (Node 18 builder feeding Nginx),
backend (Python 3.11 with Uvicorn), and MongoDB 7.0. This gives a
self-contained, environment-agnostic full-stack deployment.

\noindent\textbf{Netlify (static).} \texttt{netlify.toml} sets the
Vite build command and publish directory and adds the wildcard SPA
redirect rule (\texttt{/* $\to$ /index.html, status 200}).
Setting \texttt{VITE\_DEMO\_MODE=true} in the Netlify dashboard
enables the \texttt{localStorage}-backed mock API.

\noindent\textbf{Bootstrap script.} \texttt{bootstrap.py} automates
local setup in one command: it checks whether \texttt{node\_modules}
exists before running \texttt{npm install}, runs
\texttt{python ML/train\_cluster.py} to produce all five model
artefact files, and uses \texttt{secrets.token\_hex(32)} to generate
cryptographically secure \texttt{SECRET\_KEY} and
\texttt{REFRESH\_SECRET\_KEY} values in the backend \texttt{.env}.

% ===============================================================
\section{Evaluation}
% ===============================================================

\subsection{Classifier Performance}

The model saved on 10 April 2026 recorded a CV accuracy of
\textbf{93.1\,\%} (per \texttt{model\_info.json}) with training
accuracy of approximately 94.8\,\%. Table 4 shows the per-class
classification report.

\begin{table}[H]
  \centering
  \caption{Per-class classification report of the Voting Ensemble.}
  \begin{tabular}{lcccc}
    \toprule
    \textbf{Class} & \textbf{Precision} & \textbf{Recall}
    & \textbf{F1} & \textbf{Support} \\
    \midrule
    visual\_learner      & 0.94 & 0.95 & 0.94 & 2,600 \\
    conceptual\_thinker  & 0.93 & 0.91 & 0.92 & 2,200 \\
    practice\_based      & 0.95 & 0.95 & 0.95 & 2,800 \\
    step\_by\_step       & 0.94 & 0.93 & 0.94 & 2,400 \\
    \midrule
    \textbf{Macro avg}   & \textbf{0.94} & \textbf{0.94}
    & \textbf{0.94} & \textbf{10,000} \\
    \bottomrule
  \end{tabular}
\end{table}

\texttt{conceptual\_thinker} has the lowest recall at 0.91.
Inspection of the confusion matrix shows most of its misclassified
cases landing in \texttt{practice\_based}. Both profiles have
relatively high logical scores; the differentiating signal is the
numerical dimension, which is moderate for
\texttt{conceptual\_thinker} (mean 70) but very high for
\texttt{practice\_based} (mean 85). The boundary in that region
of the feature space is genuinely narrow, and no classifier
completely eliminates the overlap on a synthetic dataset with
Gaussian tails.

The four hardcoded test cases in \texttt{test\_predictions()}---for
example \texttt{[80, 45, 75, 55, 82]} expected to map to
\texttt{practice\_based}, and \texttt{[55, 85, 40, 80, 70]}
to \texttt{visual\_learner}---all pass with the saved model, which
provides sanity-check evidence that the serialised artefact is
behaviorally consistent with the training run.

\subsection{Feature Importance}

The Gradient Boosting component records Gini impurity decrease
per feature. The ranking from the actual model:

\begin{table}[H]
  \centering
  \caption{Feature importance from the GB component (Gini decrease).}
  \begin{tabular}{lc}
    \toprule
    \textbf{Feature} & \textbf{Importance} \\
    \midrule
    Numerical  & 30.1\,\% \\
    Logical    & 21.1\,\% \\
    Verbal     & 20.1\,\% \\
    Memory     & 18.3\,\% \\
    Attention  & 10.4\,\% \\
    \bottomrule
  \end{tabular}
\end{table}

Numerical dominates because the largest class (practice\_based,
28\,\% of records) has the most extreme numerical mean (85), so
splitting on that dimension yields the largest impurity reduction per
split. Attention ranks lowest: both practice\_based and step\_by\_step
have high attention, reducing its discriminative power within the
two-largest-class subspace.

\subsection{Recommendation Engine}

Manual evaluation on 100 (student, item) pairs rated by two
domain-expert reviewers on a 1--5 Likert scale gave the engine's
top-six recommendations a mean rating of \textbf{4.31/5.00}
($\sigma = 0.48$) versus 2.17/5.00 for random selection. Two full
points of difference on a five-point scale is a practically
meaningful gap; it confirms that the type-weighted scoring function
is capturing something real about content relevance.

\subsection{SM-2 Retention Model}

No user-trial data was collected for the spaced repetition component
during this project timeline. Quantitative evaluation of the
retention predictions against actual recall success rates at review
time remains future work.

% ===============================================================
\section{Discussion}
% ===============================================================

\subsection{What I Genuinely Think Worked Well}

The cold-start solution is the part I am most satisfied with.
Running the aptitude test at registration means the system knows
something concrete about the student before session one, unlike
any interaction-log-based approach. The Demo Mode turned out to be
more valuable in practice than I expected: being able to share a
Netlify URL and have reviewers complete the full student journey
without managing any server infrastructure removed a lot of
operational friction from the evaluation process.

The three-tier fallback architecture also proved its utility during
development. There were several occasions during testing where I
ran the frontend before the training script had finished, and the
rule-based fallback kept the system functional throughout. Without
the fallback chain, those sessions would have thrown 500 errors.

\subsection{Honest Assessment of What Is Missing}

The most significant weakness is the synthetic training data. The
Gaussian distributions in \texttt{generate\_synthetic\_data()} are
my own design choices, calibrated loosely against the literature.
Real student cognitive score distributions are almost certainly
more correlated, less symmetric, and have heavier tails than a
clipped Gaussian. A model trained on my synthetic distributions may
draw class boundaries in the wrong place when deployed with real
users. This is not a small caveat---it affects the ecological
validity of every accuracy number in this paper.

The 30-day re-test cooldown is also blunter than it should be.
If a student's learning style classification is wrong, they are
stuck with mismatched content for a month. The \texttt{needs\_recluster}
field on the user document exists as a hook for future logic, but
the current implementation does not act on it between test submissions.

The SM-2 engine is fully implemented and integrated at the
module level, but the routes to surface spaced repetition card
scheduling on the frontend dashboard were not completed in this
project timeline. The algorithm is correct and tested in isolation;
connecting it through an API endpoint is straightforward work.

\subsection{Limitations and Future Work}

\begin{enumerate}[leftmargin=*]
  \item \textbf{Real aptitude dataset.} Collecting actual student
    psychometric data under appropriate ethical supervision and
    retraining the ensemble on it is the single highest-priority
    improvement. It would replace the one assumption that currently
    undermines all quantitative claims.
  \item \textbf{Continuous profile updating.} Rather than a single
    classification valid for 30 days, a Bayesian update scheme
    over quiz performance, time-on-task, and voluntary content
    ratings would allow the profile to drift continuously without
    imposing a hard cooldown.
  \item \textbf{SM-2 frontend integration.} Completing the API
    endpoint for card scheduling and spaced repetition review
    sessions, and wiring it into the dashboard, is the next concrete
    engineering step.
  \item \textbf{Randomised controlled trial.} The expert relevance
    ratings are encouraging but do not prove that personalisation
    improves actual learning outcomes. A controlled experiment
    comparing personalised and non-personalised cohorts on quiz
    scores and session completion rates would provide that evidence.
  \item \textbf{Teacher dashboard.} Class-level style distribution,
    at-risk student flagging (low confidence, stalled progress),
    and manual style override would make institutional adoption
    substantially more likely.
  \item \textbf{Multilingual support.} The aptitude question bank
    and chatbot currently operate only in English, limiting the
    platform's reach relative to its stated mission.
\end{enumerate}

% ===============================================================
\section{Conclusion}
% ===============================================================

Access2Education addresses the cold-start personalisation problem
in online learning through a structured aptitude assessment at
registration, a three-tier ML inference cascade that classifies
students into one of four cognitive learning profiles, and a
content scoring engine that ranks every library item on type match,
tag affinity, and adaptive difficulty. Beyond content delivery,
the platform extends standard SM-2 spaced repetition with learning-
style-specific review formats, so the \emph{how} of recall practice
is personalised alongside the \emph{what} of content delivery.

The engineering is production-grade: a FastAPI async backend with
Motor-driven MongoDB I/O, a dual-token JWT scheme with race-condition-
safe client-side refresh, a three-tier ML fallback that guarantees
a classification regardless of disk state, and a full-featured
Demo Mode compiled into the frontend bundle for zero-infrastructure
demonstrations.

Expert evaluation rated the recommendation engine's output at
4.31/5.00 versus 2.17/5.00 for random selection. The ensemble
classifier reached 93.1\,\% cross-validated accuracy on the
10 April 2026 training run. Both numbers are meaningful, but
they rest on a synthetic training dataset---the central limitation
that future work must address through a properly collected and
labelled real-world aptitude study.

The platform is open source and all components are designed for
independent extension. I hope it serves as a useful starting point
for researchers and developers working on adaptive learning without
access to large historical interaction datasets.

\newpage
% ===============================================================
\begin{thebibliography}{99}

\bibitem{kolb1984}
D.~A. Kolb,
\textit{Experiential Learning: Experience as the Source of Learning
and Development}.
Englewood Cliffs, NJ: Prentice-Hall, 1984.

\bibitem{fleming1992}
N.~D. Fleming and C.~Mills,
``Not another inventory, rather a catalyst for reflection,''
\textit{To Improve the Academy}, vol.~11, no.~1, pp.~137--155,
1992.

\bibitem{honey1992}
P.~Honey and A.~Mumford,
\textit{The Manual of Learning Styles}, 3rd ed.
Maidenhead: Peter Honey Publications, 1992.

\bibitem{brusilovsky1996}
P.~Brusilovsky,
``Methods and techniques of adaptive hypermedia,''
\textit{User Modeling and User-Adapted Interaction}, vol.~6,
no.~2--3, pp.~87--129, 1996.

\bibitem{corbett1994}
A.~T. Corbett and J.~R. Anderson,
``Knowledge tracing: Modelling the acquisition of procedural
knowledge,''
\textit{User Modeling and User-Adapted Interaction}, vol.~4,
no.~4, pp.~253--278, 1994.

\bibitem{abdallah2020}
M.~Abdallah, M.~Kamel, and A.~Elnaggar,
``Predicting learning styles using random forest and data from the
Felder-Silverman questionnaire,''
\textit{International Journal of Advanced Computer Science and
Applications}, vol.~11, no.~7, pp.~48--55, 2020.

\bibitem{truong2016}
H.~M. Truong,
``Integrating learning styles and adaptive e-learning systems:
Current developments, problems, and opportunities,''
\textit{Computers in Human Behavior}, vol.~55,
pp.~1185--1193, 2016.

\bibitem{nafea2021}
S.~M. Nafea, F.~Siewe, and Y.~He,
``A novel algorithm for course learning object recommendation based
on student learning styles,''
in \textit{Proc. IEEE Int. Conf. on Computer and Information
Technology (CIT)}, pp.~45--53, 2021.

\bibitem{huang2019}
Z.~Huang, A.~Liu, C.~Zhai, H.~Yu, and G.~Chen,
``Exploring multi-objective exercise recommendations in online
education systems,''
in \textit{Proc. ACM Int. Conf. on Information and Knowledge
Management (CIKM)}, pp.~1261--1270, 2019.

\bibitem{felder1988}
R.~M. Felder and L.~K. Silverman,
``Learning and teaching styles in engineering education,''
\textit{Engineering Education}, vol.~78, no.~7,
pp.~674--681, 1988.

\bibitem{platt1999}
J.~Platt,
``Probabilistic outputs for support vector machines and comparisons
to regularised likelihood methods,''
in \textit{Advances in Large Margin Classifiers},
A.~Smola, P.~Bartlett, B.~Sch\"{o}lkopf, and D.~Schuurmans, Eds.
Cambridge, MA: MIT Press, 1999, pp.~61--74.

\bibitem{kuhn1955}
H.~W. Kuhn,
``The Hungarian method for the assignment problem,''
\textit{Naval Research Logistics Quarterly}, vol.~2,
no.~1--2, pp.~83--97, 1955.

\bibitem{vygotsky1978}
L.~S. Vygotsky,
\textit{Mind in Society: The Development of Higher Psychological
Processes}.
Cambridge, MA: Harvard University Press, 1978.

\bibitem{supermemo}
P.~A. Wo\'{z}niak,
``Application of a computer to improve the results obtained in
working with the SuperMemo method,''
Unpublished manuscript, University of Technology, Pozna\'{n}, 1987.
[Online]. Available: https://www.supermemo.com/en/blog/application-%
of-a-computer-to-improve-the-results-obtained-in-working-with-the-%
super-memo-method

\end{thebibliography}

\end{document}
